%% kapa: exam=1 round=30 year=2019 period=1 form=A subject=real_estate subject_ko=부동산학원론 seq=1 number=81
%% source: Q-Net 감정평가사 기출문제 (cst00309) artlSeq=5205205 file=extracted/1008/5205205/2019년 제30회 감정평가사 1차시험 문제지/2019년 제30회 감정평가사 1차시험 문제지/2019년도 제30회 감정평가사 1차 1교시 A형.hwp
%% conversion: kordoc 4.24.51 -> tools/gen1.py (auto); review=proofread-sonnet; flags=none
\begin{question}{81}
토지의 특성에 관한 설명으로 옳지 않은 것은?
\begin{choices}
  \item 부동성은 부동산 활동 및 현상을 국지화하여 지역특성을 갖도록 한다.
  \item 부증성은 생산요소를 투입하여도 토지 자체의 양을 늘릴 수 없는 특성이다.
  \item 영속성은 토지관리의 필요성을 높여 감정평가에서 원가방식의 이론적 근거가 된다.
  \item 개별성은 대상토지와 다른 토지의 비교를 어렵게 하며 시장에서 상품 간 대체관계를 제약할 수 있다.
  \item 인접성은 물리적으로 연속되고 연결되어 있는 특성이다.
\end{choices}
\end{question}
\answer{3}
